%% kapa: exam=1 round=31 year=2020 period=1 form=A subject=civil_law subject_ko=민법 seq=15 number=15
%% source: Q-Net 감정평가사 기출문제 (cst00309) artlSeq=5207219 file=extracted/1008/5207219/A형/A형/2020년도 제31회 감정평가사 1차 시험_1교시_A형.pdf
%% conversion: kordoc 4.24.51 -> tools/gen1.py (auto); review=proofread-sonnet; flags=none
\begin{question}{15}
법률행위의 효력에 관한 설명으로 옳은 것을 모두 고른 것은? (다툼이 있으면 판례에 따름)
\begin{framedtext}
ㄱ. 매매계약을 체결하면서 양도소득세를 면탈할 의도로 소유권이전등기를 일정기간 유보하는 약정은 반사회질서행위로 볼 수 없다. \par
ㄴ. 경매목적물과 매각대금이 현저하게 공정을 잃은 경우에도 그 경매는 불공정한 법률행위에 해당하지 않는다. \par
ㄷ. 도박에 쓸 것을 알면서 빌려준 금전을 담보하기 위하여 저당권을 설정한 사람은 저당권설정등기의 말소를 청구할 수 있다. \par
\end{framedtext}
\begin{choices}
  \item ㄱ
  \item ㄴ
  \item ㄱ, ㄷ
  \item ㄴ, ㄷ
  \item ㄱ, ㄴ, ㄷ
\end{choices}
\end{question}
\answer{5}
